\documentclass{article}
\usepackage[utf8]{inputenc}
\usepackage[russian]{babel} % For Russian language support

\title{Пример научной статьи}
\author{Автор}
\date{\today}

\begin{document}

\maketitle

\section{Введение}
Настоящая статья представляет собой пример документа, написанного на языке LaTeX. 
LaTeX является мощной системой подготовки документов, особенно подходящей для научных и математических работ, 
благодаря своей превосходной системе набора формул. В этой статье мы рассмотрим основную структуру документа 
типа article и продемонстрируем использование различных элементов форматирования.

LaTeX позволяет легко создавать структурированные документы с автоматической нумерацией разделов, 
ссылками на литературу и библиографией. Это делает его популярным инструментом среди ученых, 
студентов и преподавателей по всему миру.

\section{Основные элементы структуры}
Документ типа article включает в себя титульный лист, разделы, подразделы, списки, формулы и другие элементы. 
Каждый из этих элементов может быть легко оформлен с помощью соответствующих команд LaTeX.

\section{Заключение}
В заключение отметим, что LaTeX остается одним из наиболее мощных и гибких инструментов для подготовки 
научных документов. Его использование позволяет сосредоточиться на содержании работы, 
в то время как форматирование осуществляется автоматически в соответствии с заданными стилями. 
Это делает LaTeX незаменимым инструментом для написания статей, диссертаций и книг.

\end{document}